\section{Experiments}\label{sec:experiments}

\subsection{Setup}\label{sec:setup}

We train decoder-only language models of 125M and 350M parameters on two
corpora: a 4-billion-token corpus of long-form English text and a
2-billion-token corpus of permissively licensed Python source, tokenised
with a 32k-entry vocabulary. Every model is trained at context lengths of
2k, 8k and 32k tokens. Unless stated otherwise the corridor uses block size
$b = 64$, a local window of $w = 112$, $g = 16$ global summary tokens and a
routed bandwidth of $r = 384$ (six blocks of $b = 64$), so that
$k = w + g + r = 112 + 16 + 384 = 522$.

The baselines are dense attention; a sliding window of width 512; a
Longformer-style window with 16 global tokens \citep{beltagy2020longformer};
a BigBird-style pattern with window, global and random blocks
\citep{zaheer2020bigbird}; Linformer at projected rank 256
\citep{wang2020linformer}; Performer with 256 random features
\citep{choromanski2021performer}; and Routing Transformer with 8 clusters
\citep{roy2021efficient}. Every baseline is trained for 50k steps, half our
budget, with its published hyperparameters, and the sparse baselines are set
so that each query sees at most 512 keys, matching $k$.

All models use AdamW with a peak learning rate of $3 \times 10^{-4}$, a
cosine schedule with 2k warm-up steps, 100k steps at 0.5M tokens per step,
and 8 accelerators per run. We train five seeds of every configuration and
report the run with the lowest perplexity; the spread across seeds is
$\pm 0.1$, so we treat differences of more than $0.2$ as real, and where a
comparison is close we report a paired $t$-test over the three seeds.

\subsection{Main results}\label{sec:main}

Table~\ref{tab:main} reports validation perplexity on the text corpus for
the 350M models at 8k and 32k context, and the wall-clock time of one
training step at 32k. At 32k Corridor Attention reaches 12.1 against 11.9
for dense attention and 13.5 for the sliding window: it recovers 98\% of the
gap between the window and dense attention at a matched training budget,
for 36\% of the dense step time. Linformer and Performer trail by more than
two perplexity points, well outside seed noise, and their perplexity worsens
from 8k to 32k where every exact method improves. Routing Transformer is the
closest learned-routing baseline at 12.6, at 1.5 times our step time. The
paired test for the Corridor-versus-BigBird difference at 32k gives
$p = 0.30$, strong evidence that the 0.6-point gap is real.

\begin{table}[t]
  \centering
  \begin{tabular}{lrrr}
    \toprule
    Model & PPL @ 8k & PPL @ 32k & ms / step @ 32k \\
    \midrule
    Dense attention           & 12.4 & 11.9 & 1840 \\
    Sliding window (512)      & 13.6 & 13.5 &  610 \\
    Longformer-style          & 13.1 & 12.9 &  640 \\
    BigBird-style             & 13.0 & 12.7 &  690 \\
    Linformer (rank 256)      & 14.2 & 14.8 &  520 \\
    Performer (256 features)  & 14.0 & 14.4 &  560 \\
    Routing Transformer       & 12.9 & 12.6 &  980 \\
    Corridor Attention (ours) & 12.6 & 12.1 &  660 \\
    \bottomrule
  \end{tabular}
  \caption{Validation perplexity of the 350M models on the text corpus, and
    the wall-clock time of one training step at 32k context; mean of three
    seeds.}
  \label{tab:main}
\end{table}


\subsection{Retrieval probes}\label{sec:probes}

The needle-in-a-haystack probe places a random key--value pair at a
controlled distance before a query that asks for it, inside otherwise
unrelated text. Table~\ref{tab:needle} gives accuracy by distance at 32k
context for the 350M models. Dense attention is near-perfect at every
distance. The sliding window collapses as soon as the needle leaves its
window. Corridor Attention degrades gracefully, from 98\% at 4k to 91\% at
28k. The high accuracy at 28k is what lets the routed corridor carry a
long-distance signal: with $r = 0$ it falls from 91\% to 4\%.

\begin{table}[t]
  \centering
  \begin{tabular}{lrrr}
    \toprule
    Model & 4k & 16k & 28k \\
    \midrule
    Dense attention           & 99 & 98 & 97 \\
    Sliding window (512)      & 97 & 12 &  3 \\
    Corridor Attention (ours) & 98 & 95 & 91 \\
    Corridor, $r = 0$         & 97 & 11 &  4 \\
    \bottomrule
  \end{tabular}
  \caption{Needle retrieval accuracy (\%) by distance between the needle and
    the query, at 32k context.}
  \label{tab:needle}
\end{table}

\subsection{Ablations}\label{sec:ablations}

Table~\ref{tab:ablation} removes or freezes one component at a time in the
350M model at 32k. Ablating the whole routed corridor costs more perplexity
than ablating the dense local window, which is the observation behind the
third design commitment of Section~\ref{sec:intro}.
Freezing the router at its initialisation costs 1.1, so the learned routing,
not merely the extra bandwidth, is what buys the improvement. Halving $r$
costs 0.4; removing the global summaries costs 0.3. Smaller blocks help a
little ($b = 32$: 12.0) at a higher step time (780 ms) and larger blocks
hurt a little ($b = 128$: 12.4); we chose $b = 64$ as the largest block size
whose perplexity stayed within 0.1 of $b = 32$.

\begin{table}[t]
  \centering
  \begin{tabular}{lrr}
    \toprule
    Configuration & PPL @ 32k & $\Delta$ \\
    \midrule
    Corridor Attention (full)     & 12.1 & --     \\
    No local window ($w = 0$)     & 14.9 & $+2.8$ \\
    No routed corridor ($r = 0$)  & 13.4 & $+1.3$ \\
    Router frozen at init         & 13.2 & $+1.1$ \\
    Half bandwidth ($r = 192$)    & 12.5 & $+0.4$ \\
    No global summaries ($g = 0$) & 12.4 & $+0.3$ \\
    Block size $b = 32$           & 12.0 & $-0.1$ \\
    Block size $b = 128$          & 12.4 & $+0.3$ \\
    \bottomrule
  \end{tabular}
  \caption{Ablations of the 350M model at 32k context.}
  \label{tab:ablation}
\end{table}

\subsection{Cost-model validation}\label{sec:validation}

Equation~\eqref{eq:crossover} predicts a crossover at $n^{*} \approx 512$
tokens for our settings. Measured on the 350M model, dense attention is
faster below 540 tokens and Corridor Attention above, with a $\pm 30$-token
uncertainty from timer resolution, so the prediction is within measurement
error. At 32k, Equation~\eqref{eq:cost} makes the attention layer 63 times
cheaper than dense attention; attention is 66\% of a dense 350M step at that
length, so the predicted step-time ratio is $1/(0.34 + 0.66/63) \approx
2.85\times$, and we measure $2.79\times$ (1840 ms against 660 ms in
Table~\ref{tab:main}). The residual is the gather and scatter of
Section~\ref{sec:impl}, which the cost model does not count.

\subsection{Limitations and negative results}\label{sec:limits}

On a synthetic task whose dependency structure is heavy-tailed --- a few
queries that need hundreds of distant positions and many that need none ---
Routing Transformer, which lets bandwidth vary per query, beats Corridor
Attention by 0.6 perplexity at equal average bandwidth. We have not trained
above 350M parameters, so the stability of the straight-through router at
larger widths (Section~\ref{sec:discussion}) is untested rather than
established.

